\chapter{Conclusion and Future Work}

\section{Conclusion, Limitations, and Future Scope}
This project report presented a complete conceptual and analytical framework for customer churn prediction in a FinTech platform using machine learning. The study began by identifying the business relevance of churn and showing why retention is a central concern for digital financial services. Unlike simple descriptive reporting, a predictive churn system attempts to estimate which customers are likely to disengage before the actual churn event takes place. This creates direct practical value because customer success teams, support teams, and retention managers can act earlier and more selectively. The report demonstrated that churn prediction is not merely a classification problem in a narrow technical sense, but a broader decision-support problem involving data preparation, feature design, model comparison, evaluation, interpretation, and communication of findings to business users.

The work made several important contributions within the scope of an academic major project. First, it surveyed relevant literature and showed that the field has moved toward ensemble-based methods, imbalance-aware modeling, and more deployment-conscious system design. Second, it translated those research insights into a proposed system architecture that includes data ingestion, preprocessing, feature engineering, model training, prediction service, and dashboard presentation. Third, it identified a realistic feature schema for a FinTech churn dataset and described how raw variables can be transformed into more meaningful indicators such as transaction intensity and complaint frequency. Fourth, it compared four relevant algorithms, namely Logistic Regression, Decision Tree, Random Forest, and XGBoost, and explained why ensemble methods are especially suitable for structured customer data. Fifth, it established a multi-metric evaluation framework that gives due importance to recall, F1-score, and ROC-AUC rather than relying only on accuracy.

Among the compared models, XGBoost emerged as the strongest candidate in the illustrative evaluation. This result is consistent with current churn prediction literature, which frequently reports that boosted tree methods perform well on tabular business datasets \cite{imani2024,bhattacharjee2026}. The reason is that churn behavior is shaped by interactions among multiple features rather than by one simple linear factor. Customers may leave because of low activity combined with low tenure, or because of repeated complaints combined with moderate usage decline, or because of a mismatch between subscription type and service expectations. Ensemble models are better able to capture these layered relationships. At the same time, the report also acknowledged the importance of interpretability and noted that techniques such as feature importance analysis, SHAP, and LIME can make strong models more transparent to business stakeholders.

Despite the strengths of the proposed system, the study has several limitations. The first limitation is that the dataset is assumed rather than collected from a live institutional FinTech source. Although the schema and metrics are grounded in published literature, actual organizational data may contain additional complexities such as temporal drift, missing behavioral context, or domain-specific variables that are not included here. The second limitation is that the performance results are illustrative and intended to represent expected behavior rather than final experimental proof. These results should be replaced with actual model outputs when the full implementation is run on a local machine with proper data access. The third limitation is that the project stops short of real-time deployment and controlled business experimentation. In practice, a churn prediction system is only fully validated when it is connected to retention actions and measured through campaign outcomes.

Future scope naturally follows from these limitations. One important extension is the integration of real-time or periodic batch scoring into a production environment so that customer risk is updated continuously. Another is the addition of explainable AI methods that provide local and global reasoning behind predictions, improving trust among non-technical decision-makers. Further work may also include automated retraining schedules to handle concept drift, richer feature sources such as support transcripts or customer sentiment signals, and threshold optimization based on business cost analysis. A particularly valuable next step would be A/B testing, where retention interventions guided by the churn model are compared against standard business practice. Such experimentation would allow the institution or company to measure not only predictive accuracy, but true business impact. In summary, this report establishes a strong foundation for a practical churn prediction system and shows how machine learning can support more intelligent, timely, and efficient retention strategies in FinTech environments.
